\documentclass{homework}
% \usepackage{graphicx} % Required for inserting images
\usepackage{amsmath}
\usepackage{gensymb}
\usepackage{subfig}
\usepackage{float}

\class{ECEN 5407}
\title{Homework 2}
\author{Sean Jennings}
\date{September 20, 2023}

\begin{document}
\maketitle

\section{Problem 1}
\begin{letterenum}
\item The instantaneous power is given by the equation
\begin{equation}
    P = \frac{1}{2} \rho A U^3 C_P
\end{equation}
where $\rho$ is the air density, $A$ is the swept area,
$U$ is the wind speed, and $C_P$ is the power coefficient.
We assume the air is at STP and therefore has a density of
$1.225$ kg/m$^3$. This seems like a reasonable assumption because Austin
is at a low elevation of ~500 feet. We compute swept area by first converting
the diameter $D$ to meters
\begin{equation}
    A = \frac{\pi}{4} D^2 = 2339 \text{m/s}
\end{equation}
We also convert the wind speed to m/s before performing the computation
in Python.  Table (\ref{table:powers}) shows the instantaneous power as a function of
time.
\begin{table}[h!]
\centering
\begin{tabular}{ |c|c|}
\hline
Time & Power (kW) \\
\hline
0000 & 0.36 \\
0100 & 0.79 \\
0200 & 0.04 \\
0300 & 0.1 \\
0400 & 0.15 \\
0500 & 0.1 \\
0600 & 1.09 \\
0700 & 4.36 \\
0800 & 12.88 \\
0900 & 28.49 \\
1000 & 22.93 \\
1100 & 33.76 \\
1200 & 37.21 \\
1300 & 37.21 \\
1400 & 22.09 \\
1500 & 12.3 \\
1600 & 15.36 \\
1700 & 30.53 \\
1800 & 19.2 \\
1900 & 4.08 \\
2000 & 0.88 \\
2100 & 1.61 \\
2200 & 4.95 \\
2300 & 8.74 \\
\hline
\end{tabular}
\caption{Instantaneous Power over time}
\label{table:powers}
\end{table}

\item Figure (\ref{fig:instantaneous-power}) plots the instantaneous power
curve. As expected, it follows the shape of the wind speed curve. However,
the slopes of the curve are magnified due to the cubic relation with wind speed.

\begin{center}
    \includegraphics[width=0.8\linewidth]{WIND_SPEED.png}
    \captionof{figure}{Instantaneous Power in Austin}
    \label{fig:instantaneous-power}
\end{center}

\end{letterenum}

\section{Problem 2}
\begin{letterenum}
\item The heat rate, HR, is the amount of energy in the fuel required to generate
1 kWh of electricity. To find the heat rate, we divide the amount of
fuel energy by the total electricity generated.
\begin{equation}
    \text{HR} = \frac{75,444,200 \text{MMBtu} \cdot 10^6 \text{ Btu/MMBtu}}
        {7,561,850 \text{ MWh} \cdot 1000 \text{ kWh/MWh}} = 9977.  \text{ Btu (/kWh)}
\end{equation}

To compute the efficiency as a percentage, we divide the number of BTU
in a kWh by the heat rate. 1 kWh is equivalent to 3,412 Btu, so the
efficiency is
\begin{equation}
    \eta = \frac{3412 \text{ Btu / kWh}}{9977 \text{ Btu / kWh}} = 34.2 \%
\end{equation}

\item Figures (\ref{fig:heat-input}) and (\ref{fig:heat-rate}) show the
heat input and heat rate curves, respectively. We assume that each row of the
table represents an hour of time, so that the energy generation in MWh of the
left column is equivalent to the power in MW which is displayed. This
assumption results in heat rate values of a reasonable magnitude. E.g.
the heat rates correspond to a mean efficiency of 37\%.

The author guesses this is a coal plant since it is running at a nearly
constant rate: coal plants typically have slow ramp times, high minimum generation
factors, and long start-up times. Coal plants at traditional combustion
pressures have effiencies ~35-38\% so the efficiency also falls in range.

The slow ramping could also make it a candidate for a nuclear
plant, but the capacity of ~360MW is less than the 1GW typical of a nuclear
reactor unit -- at least until we have some SMRs (small modular reactors)
come online!

\begin{center}
    \includegraphics[width=0.6\linewidth]{heat_input_curve.png}
    \captionof{figure}{Heat input curve}
    \label{fig:heat-input}
\end{center}

\begin{center}
    \includegraphics[width=0.6\linewidth]{heat_rate_curve.png}
    \captionof{figure}{Heat rate curve}
    \label{fig:heat-rate}
\end{center}
\end{letterenum}

\section{Problem 3}
We will assume that the current of 3A which is displayed in the diagram
represents the RMS current of the AC circuit. The complex power is given by:
\begin{equation}
    S = VI^* = |I|^2 R = (3\text{ A})^2 * 100 \ohm = 900 \text{W} + j0 \text{ var}
\end{equation}
Since the circuit is purely resistive, the phase angle between voltage and current
is zero. The real component of the complex power, $P$, is real power, and the imaginary
component, $Q$, is reactive power.
\begin{equation}
    S = P + jQ
\end{equation}
The real power is $\boxed{$900W$}$.

\noindent The reactive power is $\boxed{$0var$}$.

\noindent Apparent power is the magnitude of complex power, with units VA by convention,
and is equal to $\boxed{$900VA$}$. It is equal to the real power since the load is purely
resistive.

\noindent The power factor is the cosine of the phase angle, so in this case: $\cos(0) = 1$.


If a capacitor were added to the circuit, the current would lead the voltage,
i.e. would have a positive phase angle.
Since complex power multiplies $V$ by the conjugate of $I$, so the reactive
power absorbed would be negative. The power factor would be less than 1.

\section{Problem 4}

The total complex power delivered by the source is the sum of the power
absorbed by the loads:
\begin{equation}
    S = \sum_i S_i = S_R + S_C + S_L = 20 \text{W} + j(-5 \text{ VAr})
        + j(20 \text{ VAr}) = 20 \text{ W} + j15 \text{ VAr}
\end{equation}
So the apparent power is $|S| = \sqrt{20^2 + 15^2} = 25$ VA.

The power factor is
\begin{equation}
    \text{pf} = \tan^-1(Q / P) = \tan^-1(15/20) \approx 0.64
\end{equation}
This power factor is quite low. A lower power
factor means that the power system is not running at peak efficiency, since
a significant amount of the total apparent power generated is not being
consumed as real power. Since customer loads such as motors are usually
inductive, a distribution company will install shunt capacitors that will 
deliver reactive power locally, so that reactive power does not need to
be transmitted over distribution lines. This minimizes the line losses
associated with the higher currents that would be required to deliver reactive
power. Typically, a power factor under 0.85 in a distribution system 
is considered poor.

Figure (\ref{fig:power-tri-component}) plots the complex power vector
decomposed into the contributions by each ciruit component.
Figure (\ref{fig:power-tri-real}) plots the
complex power vector decomposed into real and reactive parts.

\begin{center}
    \includegraphics[width=0.6\linewidth]{power_triangle_component.png}
    \captionof{figure}{Power triangle: by component}
    \label{fig:power-tri-component}
\end{center}

\begin{center}
    \includegraphics[width=0.6\linewidth]{power_triangle_real.png}
    \captionof{figure}{Power triangle: real vs reactive}
    \label{fig:power-tri-real}
\end{center}

\section{Problem 5}

\begin{letterenum}
    \item The base voltage levels for this power system are 161 kV and 69kV.
    This is an example of a sub-transmission system. Distribution system 
    voltage range from 4-35kV. Long-distance transmission lines has
    voltages above 200kV to reduce line losses. Subtransmission lines
    have voltages less than 200kV, and carry power between transmission
    and distribution substations. The voltages in this diagram are
    in this intermediate range.

    \item Table (\ref{table:power-by-bus}) shows the the real, reactive, and apparent power for each
    corresponding to each bus with a load. For each bus $i$, apparent power
    consumed is
    \begin{equation}
        |S_i| = \sqrt{P_i^2 + Q_i^2},
    \end{equation}
    where $i\in B$ and $B$ is the set of each bus which has a load.

    Then, the total real and reactive power for the full system is given by
    the sum of the $P$ and $Q$ columns of table (\ref{table:power-by-bus}),
    respectively:
    \begin{equation}
        P_{total} = \sum_{i \in B} P_i = 585 \text{ MW} \qquad
        Q_{total} = \sum_{i \in B} Q_i = 145 \text{ MVAr}
    \end{equation}
    We see the largest load (by every measure) occurs at bus 5.

    \begin{table}[h!]
    \centering
    \begin{tabular}{|c|c|c|c|}
        \hline
        Bus & $P$ (MW) & $Q$ (MVAr) & $|S|$ (VA) \\
        \hline
        4 & 60 & 10  & 60.83  \\
        5 & 100 & 30 & 104.4  \\
        6 & 80 & 15  & 81.39  \\
        7 & 90 & 20  & 92.2   \\
        8 & 40 & 5   & 40.31  \\
        9 & 10 & 5   & 11.18  \\
        10 & 15 & 10 & 18.03  \\
        11 & 75 & 15 & 76.49  \\
        12 & 40 & 15 & 42.72  \\
        13 & 30 & 10 & 31.62  \\
        14 & 35 & 10 & 36.4   \\
        15 & 10 & 0  & 10.0   \\
        \hline
    \end{tabular}
    \caption{Real, Reactive, and Apparent Power by Bus}
    \label{table:power-by-bus}
    \end{table}
    
















    
\end{letterenum}









\end{document}
